\section{Testing \& infrastructure}\label{sec:testing}

This section describes the infrastructure a new engineer inherits at the pinned commit
\code{254cd41}: six GoogleTest binaries registered with ctest as 246 entries, a small
family of oracle patterns that the suites use instead of golden values, a committed
fixture corpus the tests locate through compile-time absolute paths, a Clang-only
\code{llvm-cov} coverage pipeline living in its own build tree, and one continuous
integration (CI) workflow that drives five platforms through the same CMake
presets developers use locally (\cref{sec:overview} introduced the presets and the
warning policy; this section covers the rest). The infrastructure is disciplined, but it
has sharp edges, and each is marked where it lives: a coverage target that predates the
newest test suite, one stale comment claiming tooling that does not exist, and an install
rule that would ship the vendored dependencies alongside the application.

\subsection{Six binaries, two registrations}\label{sec:testing-binaries}

Test sources live next to what they test: \srcf{model/tests/} and \srcf{sim/tests/} hold
the per-layer unit suites, and the top-level \srcf{tests/} directory holds everything
that crosses layers --- integration, propagator behavior, and the two suites that drive
the command-line interface (CLI). \Cref{tab:testing-binaries} lists the six binaries with
their live per-case counts.

\begin{table}[htbp]
\centering\small
\begin{tabularx}{\linewidth}{@{}lr>{\raggedright\arraybackslash}Xl@{}}
\toprule
Binary & Cases & What it exercises & Defined at \\
\midrule
\code{model\_tests} & 144 & part composition and inertia, motors and thrust curves (with the RASP parser), the concrete part types, Barrowman aero, the part factory, placement types, the resolver sweep, the diagnostics gate & \src{model/tests/CMakeLists.txt}{2} \\
\addlinespace
\code{sim\_tests} & 27 & the atmosphere models, \code{Environment}, the RK45 solver, spherical geoid and gravity, integrator selection & \src{sim/tests/CMakeLists.txt}{2} \\
\addlinespace
\code{integration\_tests} & 43 & end-to-end physics through the controller, motor-database and design persistence, the placement-invariance gate, the poke-through regression & \src{tests/CMakeLists.txt}{4} \\
\addlinespace
\code{propagator\_tests} & 6 & the propagator's termination guards and trajectory statistics, driven by a mock \code{Propagatable} & \src{tests/CMakeLists.txt}{32} \\
\addlinespace
\code{cli\_tests} & 3 & a scripted REPL build $\to$ save $\to$ reload $\to$ fly session & \src{tests/CMakeLists.txt}{48} \\
\addlinespace
\code{design\_matrix\_tests} & 16 & the committed fixture matrix, round-trip fidelity, and the 1/4A $\to$ M flight ladder & \src{tests/CMakeLists.txt}{69} \\
\bottomrule
\end{tabularx}
\caption{The six test binaries at the pinned commit, with per-case counts taken live from
each binary's \code{-{}-gtest\_list\_tests} output ($144+27+43+6+3+16 = 239$ cases across
49 distinct gtest suites).}\label{tab:testing-binaries}
\end{table}

The link lines mirror the layering of \cref{sec:overview}: each binary links exactly the
topmost layer it exercises. \code{model\_tests} links \code{model} and \code{utils}
\src{model/tests/CMakeLists.txt}{22}; \code{sim\_tests} links \code{sim}
\src{sim/tests/CMakeLists.txt}{12}; \code{propagator\_tests} links \code{model} --- which
transitively pulls in \code{sim} and \code{utils} and supplies the out-of-line
\code{Propagatable} constructor symbol \src{tests/CMakeLists.txt}{37};
\code{integration\_tests} links \code{qtrocket\_core}, the same controller archive the
executables link \src{tests/CMakeLists.txt}{19}; and the two CLI-driving suites link the
\code{cli} library, so they exercise the identical \code{Repl} object code users type
into \src{tests/CMakeLists.txt}{56}, \src{tests/CMakeLists.txt}{78}.

Every binary is registered with ctest twice. First, per case:
\cppinline|gtest_discover_tests| runs each binary at build time and registers every
\code{Suite.Case} as its own ctest entry \src{model/tests/CMakeLists.txt}{29},
\src{sim/tests/CMakeLists.txt}{18}, \src{tests/CMakeLists.txt}{25} --- 239 entries on
this tree, which is what makes per-case filtering like \code{ctest -R PartTests} work.
Second, in aggregate: one \cppinline|add_test(NAME qtrocket_<suite>_tests ...)| per
binary runs the whole executable as a single ctest entry
\src{model/tests/CMakeLists.txt}{31}, \src{sim/tests/CMakeLists.txt}{20},
\src{tests/CMakeLists.txt}{27}. Seven aggregates exist for six binaries --- the
design-matrix binary registers a second aggregate, as \cref{sec:testing-heavy} explains --- for a
total of $239 + 7 = 246$ registered tests, counted live with \code{ctest -N} at the
pinned commit.

\begin{designrule}[Aggregate names are the selection contract]
Both CI and the coverage pipeline run \emph{only} the aggregates: CI filters with
\cppinline|ctest -R 'qtrocket_*'|
\src{.github/workflows/cmake-multi-platform.yml}{74} and the coverage script defaults to
\cppinline|^qtrocket_.*tests$| \src{cmake/RunLlvmCoverage.cmake}{14}. Note that
\cppinline|'qtrocket_*'| is a regular expression, not a glob --- it means ``\code{qtrocket}
followed by zero or more underscores, anywhere in the name'' --- and it selects exactly
the seven aggregates only because no discovered \code{Suite.Case} name contains the
substring \code{qtrocket}. The corollary binds new work: a future suite that forgets to
register a \code{qtrocket\_}-prefixed aggregate silently drops out of both CI and
coverage, with no error anywhere.
\end{designrule}

The net has known holes, stated here and inventoried in \cref{sec:incomplete}: no test
binary covers \srcf{gui/} or \srcf{visualizer/} (the visualizer is also the only project
target without a \code{qtrocket\_enable\_coverage} call ---
\srcf{visualizer/CMakeLists.txt} contains none, where every other target has one); and
\srcf{utils/} has no test directory of its own \srcf{utils/CMakeLists.txt}, so
\code{CurlConnection} is never executed by any suite. Each of these is an absence
verified by searching the tree, not an impression.

\subsection{Oracle patterns}\label{sec:testing-oracles}

The suites share a house style: assertions state \emph{properties that must hold} ---
closed-form values derived independently, invariants, inequalities --- rather than golden
numbers snapshotted from the code under test. Four patterns recur, and each is worth
learning by example because new tests are expected to follow them.

\begin{description}

\item[Independent numeric oracles.] The inertia-tensor suite states the rule in its
header: the tests pin the per-unit-mass helpers ``against INDEPENDENT numeric oracles
(disk / surface / volume-mesh integration) --- never against the closed form they
validate'' \src{model/tests/InertiaTensorsTests.cpp}{8}. The solid cone is checked
against a Riemann stack of 400{,}000 thin disks
\src{model/tests/InertiaTensorsTests.cpp}{28}, the conical shell against a stack of hoops
\src{model/tests/InertiaTensorsTests.cpp}{55}, and the trapezoidal fin set against a
brute-force volume mesh over $N$ prisms
\src{model/tests/InertiaTensorsTests.cpp}{92} --- the stated acceptance gate for the fin
set's closed-form algebra. Agreement is demanded at $10^{-5}$ relative for the cone
\src{model/tests/InertiaTensorsTests.cpp}{157} and $0.5\%$ for the mesh
\src{model/tests/InertiaTensorsTests.cpp}{209}; the textbook closed forms are then
additionally pinned exactly, e.g. $I_{zz} = \tfrac{3}{10}R^2$ per unit mass at
$10^{-15}$ \src{model/tests/InertiaTensorsTests.cpp}{161}. The logic is the point: a
closed form checked against its own algebra proves nothing, but an independent
discretization of the defining integral converges to the same number only if both are
right. (The tensors themselves are the subject of \cref{sec:parts}.)

\item[Literal-anchor regression pins.] Where a bug was found and fixed, its exact
numbers are frozen into an assertion with the derivation written beside it. The
poke-through regression loads the one deliberately self-intersecting fixture through the
real file reader and pins the diagnostic's location and depth:
\cppinline|EXPECT_NEAR(d.zWorld, -0.26, 1e-9)| and penetration
$0.0376 - 0.0395\,(0.26/0.30) \approx 0.003367\m$
\src{tests/PokeThroughRegressionTests.cpp}{63}, then asserts that the composite gate
refuses the failed solve --- \cppinline|EXPECT_THROW| on \code{getCompositeI}
\src{tests/PokeThroughRegressionTests.cpp}{67} --- the fail-closed posture of
\cref{sec:placement}, exercised end to end. The same file pins every resolved station of
the whitepaper's worked example to hand-computed literals
\src{tests/PokeThroughRegressionTests.cpp}{70}. The invariance suite carries two more:
the solid nose cone's local center of mass (CM) hand-pinned to $-0.225\m$ at
$10^{-12}$, with the $\bar{h} - L = -3L/4$ derivation in the comment
\src{tests/PlacementInvarianceTests.cpp}{445}, and the fin-set CM recomputed from the
fixture's raw trapezoid geometry rather than read back through the part
\src{tests/PlacementInvarianceTests.cpp}{458}. Each pin defends itself against fixture
drift --- a guard assertion such as ``fixture cone length changed --- re-derive the pin''
fails first if someone edits the geometry \src{tests/PlacementInvarianceTests.cpp}{450}.

\item[Invariance gates.] The placement migration (\cref{sec:placement}) rests on the
claim that it changed no physics, and the gate makes that claim falsifiable. A frozen
baseline snapshots composite mass, CM, inertia, and every part's resolved station for the
24-fixture corpus \src{tests/PlacementInvarianceTests.cpp}{66}, formatted at
\code{\%.17g} so every double round-trips to the identical IEEE-754 value
\src{tests/PlacementInvarianceTests.cpp}{88}. The live tree must reproduce it: mass and
inertia directly \src{tests/PlacementInvarianceTests.cpp}{266}, CM and stations after the
one known tip-datum shift \src{tests/PlacementInvarianceTests.cpp}{286},
\src{tests/PlacementInvarianceTests.cpp}{299}, all within $10^{-9}$ --- a tolerance sized
to admit the ULP-level drift of re-associated floating-point arithmetic and nothing
larger \src{tests/PlacementInvarianceTests.cpp}{151}, with a diagnostic test reporting
the worst observed drift so the tolerance can be judged against reality
\src{tests/PlacementInvarianceTests.cpp}{328}. The baseline file declares its own
epistemics: never regenerate it to make a failing migration pass, because ``a diff here
means the refactor changed the physics''
\src{tests/data/placement-invariance-baseline.txt}{4}. One quantity is deliberately
excluded: the center of pressure, because the migration \emph{corrected} a
CM-contaminated legacy value, so baseline conformance there would pin the bug
\src{tests/PlacementInvarianceTests.cpp}{318}.

\item[End-to-end flight sweeps.] The design-matrix suite flies the whole engine ---
reader, resolver, motor database, propagator, both integrators --- across seven reference
airframes and asserts physical inequalities rather than trajectories. Per motor flown:
apogee finite, above $1\m$ (``near-zero apogee implies a t=0 abort, not real flight''
\src{tests/DesignMatrixTests.cpp}{468}), below $10^{6}\m$
\src{tests/DesignMatrixTests.cpp}{469}, and strictly increasing along the motor ladder
--- \cppinline|EXPECT_GT(ap, prev)| \src{tests/DesignMatrixTests.cpp}{470}. The
atmosphere sweep flies each motor twice from identical state --- vacuum with
\code{setdrag 0}, then ``US Standard 1976'' with \code{setdrag 0.75}
\src{tests/DesignMatrixTests.cpp}{542} --- and requires \cppinline|EXPECT_LT(drag, vac)|
\src{tests/DesignMatrixTests.cpp}{552}: drag dissipates energy, so apogee under drag must
be strictly lower. The two integrators of \cref{sec:simcore} must agree in vacuum to
$2\%$: \cppinline|EXPECT_NEAR(rk45, rk4, 0.02 * rk4)|
\src{tests/DesignMatrixTests.cpp}{599}. Launch-pad behavior is pinned from both sides:
the smallest motor in the corpus must lift the lightest airframe nominally
\src{tests/DesignMatrixTests.cpp}{610}, while a $0.59$~N\,s motor under a
$2000\unit{kg/m^3}$ tube must produce a ``no liftoff'' warning and an apogee below
$1\m$ \src{tests/DesignMatrixTests.cpp}{630}. A units regression rides along:
online-sourced motor mass must arrive in kilograms, not grams
\src{tests/DesignMatrixTests.cpp}{269} --- the fix for a real bug in which a
$6.1$-gram motor loaded as $6.1$ kilograms (\cref{sec:motors}). These oracles survive
legitimate physics improvements --- a better drag model still satisfies every
inequality --- but a sign error, a unit slip, or a broken ignition path fails them
immediately.

\end{description}

\begin{physbox}[Impulse classes and the motor ladder]
Model-rocket motors are graded by total impulse $I_{\mathrm{tot}} = \int F(t)\,dt$
(newton-seconds) into letter classes, each class doubling the ceiling of the one before:
A up to $2.5\unit{N\,s}$, B up to $5$, C up to $10$, and so on up to M at
$10240\unit{N\,s}$; fractional classes 1/4A and 1/2A sit below A. This is the standard
NAR/TRA convention --- domain background, not a table encoded in the repository. A
``ladder'' is a set of motors of one diameter arranged in ascending total impulse. For a
fixed airframe in vacuum under fixed gravity, more delivered impulse means more kinetic
energy at burnout and therefore a higher apogee, which is why apogee monotonicity along
the ladder is a coarse but robust integral check on the entire
motor $\to$ force $\to$ integrator chain.
\end{physbox}

\subsection{The light/heavy registration split}\label{sec:testing-heavy}

The ladder sweeps are expensive by nature: they ``fly every motor of each class's full
1/4A $\to$ M ladder, in vacuum and through the atmosphere --- with high-impulse motors
coasting to tens of km this dominates the runtime'' \src{tests/CMakeLists.txt}{85}. The
build resolves the tension between coverage and iteration speed by registering the same
binary a third time (\cref{lst:testing-heavy}): a default (``light'') aggregate that runs
all 16 cases, and a ``heavy'' aggregate that re-runs only the sweep cases --- the
\cppinline|--gtest_filter=FlightMatrix.*:Atmosphere.*| filter selects 3 of the 16 ---
tagged with the ctest label \code{heavy} and the environment variable
\code{QTROCKET\_FULL\_LADDER=1} \src{tests/CMakeLists.txt}{96}.

\begin{listing}[htbp]
\begin{minted}{cmake}
add_test(NAME qtrocket_design_matrix_tests COMMAND design_matrix_tests)

add_test(NAME qtrocket_design_matrix_heavy_tests
         COMMAND design_matrix_tests "--gtest_filter=FlightMatrix.*:Atmosphere.*")
set_tests_properties(qtrocket_design_matrix_heavy_tests PROPERTIES
   LABELS heavy
   ENVIRONMENT "QTROCKET_FULL_LADDER=1")
\end{minted}
\caption{The double registration of the design-matrix binary
\srccap{tests/CMakeLists.txt}{94}. The light entry runs the whole suite with one motor
per class; the heavy entry re-runs just the two sweep suites under the full-ladder
environment variable, labeled so ctest can include or exclude it by
cost.}\label{lst:testing-heavy}
\end{listing}

The switch lives in the test code, not the build (\cref{lst:testing-ladder}): each of the
seven \code{Ladder} entries pairs a reference airframe with its motors in ascending total
impulse \src{tests/DesignMatrixTests.cpp}{349}, and \code{motorsToFly} returns either the
full ladder or just its smallest rung:
\begin{equation*}
\mathtt{motorsToFly}(\mathtt{lad}) =
\begin{cases}
\mathtt{lad.motors} & \text{if \texttt{QTROCKET\_FULL\_LADDER} is set,}\\[2pt]
\{\,\mathtt{lad.motors.front()}\,\} & \text{otherwise.}
\end{cases}
\end{equation*}
The check is \cppinline|std::getenv(...) != nullptr|
\src{tests/DesignMatrixTests.cpp}{377} --- presence-only, so even
\code{QTROCKET\_FULL\_LADDER=0} enables the full ladder. The light run keeps the
coverage \emph{topology} (every diameter class, every airframe, both integrators) while
flying only one motor per class; under it the monotonicity loop degenerates to a single
motor, so only the heavy entry actually exercises the apogee-ordering oracle
\src{tests/DesignMatrixTests.cpp}{461}.

\begin{listing}[htbp]
\begin{minted}{cpp}
std::vector<std::string> motorsToFly(const Ladder& lad)
{
   if(std::getenv("QTROCKET_FULL_LADDER") != nullptr)
      return lad.motors;
   return { lad.motors.front() };
}
\end{minted}
\caption{The entire light/heavy switch \srccap{tests/DesignMatrixTests.cpp}{375}. The
build's heavy ctest entry sets the environment variable; the code reads only its
presence.}\label{lst:testing-ladder}
\end{listing}

The selection algebra falls out of the label, and the build file spells it out
\src{tests/CMakeLists.txt}{92}: \cppinline|ctest -R 'qtrocket_*'| runs both entries
(full coverage --- 7 tests); \code{ctest -LE heavy} is the fast development loop
(everything except the full-ladder sweeps --- 245 of the 246 entries); \code{ctest -L
heavy} runs only the sweeps (exactly 1 entry). Because the heavy entry is a strict
re-run subset --- three cases under a stronger environment --- the plain aggregate
filter gets full coverage without paying full-ladder cost on the thirteen cheap cases
twice.

\subsection{Test data and fixtures}\label{sec:testing-data}

Tests locate their committed data through two compile-time definitions:
\code{QTROCKET\_DATA\_DIR} points at the bundled motor data in \srcf{data/} and
\code{QTROCKET\_TEST\_DATA\_DIR} at the fixture corpus in \srcf{tests/data/}, both
expanded to absolute paths from \code{CMAKE\_SOURCE\_DIR}
\src{tests/CMakeLists.txt}{15}, \src{tests/CMakeLists.txt}{53},
\src{tests/CMakeLists.txt}{74}.

\begin{designrule}[Tests are cwd-independent; shipped binaries carry no source paths]
The rationale is stated at the definition site: ``Absolute paths so the tests run from
any cwd'' \src{tests/CMakeLists.txt}{13} --- a test binary invoked by ctest, an IDE, or
by hand from any directory finds the same fixtures. The definitions are \code{PRIVATE}
to the test targets, and a search across \srcf{model/}, \srcf{sim/}, \srcf{cli/}, and
\srcf{gui/} finds no other consumer of either macro --- so the executables users run
embed no build-machine paths, and the application starts with an empty motor database
(the bundled \srcf{data/Aerotech.rse} is test input, not runtime data;
\cref{sec:motors}).
\end{designrule}

The corpus itself is small and self-describing \srcf{tests/data/README.md}: 25
\code{.qrd} designs under \srcf{tests/data/designs/} --- the 24-fixture invariance corpus
spanning seven motor-diameter classes (13, 18, 24, 29, 38, 54, 75\,mm, each airframe
sized so its motor fits the body bore), plus \code{xl75\_multi\_pokethrough.qrd}, the
deliberately self-intersecting offender --- a motor-database fixture
\code{motors/estes\_small.qmd}, and the frozen invariance baseline. The corpus checks
itself: \code{designFiles()} enumerates the directory at runtime, excluding the
poke-through offender \src{tests/DesignMatrixTests.cpp}{329}, and a dedicated test pins
the result against the 24-name expected list --- failing with ``fixture count drifted
from the documented matrix'' if anyone adds or removes a design without updating the
manifest \src{tests/DesignMatrixTests.cpp}{394}.

\subsection{Coverage}\label{sec:testing-coverage}

Coverage is Clang/LLVM source-based coverage, deliberately isolated from normal
development. The \code{coverage-clang} preset inherits \code{debug-clang} but redirects
to a separate \code{build-coverage/} tree and sets \code{QTROCKET\_ENABLE\_COVERAGE=ON}
\src{CMakePresets.json}{32}, so instrumented objects never pollute \code{build/}. The
option is hard-gated to Clang: any other compiler hits a
\cppinline|message(FATAL_ERROR ...)| naming the reason --- the pipeline uses
\code{llvm-cov}, whose instrumentation format GCC does not emit
\src{CMakeLists.txt}{101}. GCC remains fully supported for building and testing through
the \code{debug-gcc}/\code{release-gcc} presets; it is only coverage that requires
Clang. Tool discovery tolerates distribution-versioned binary names
(\code{llvm-cov-21} down to \code{-17}) \src{CMakeLists.txt}{105}.

Instrumentation is applied per target by one helper, \code{qtrocket\_enable\_coverage}
\src{CMakeLists.txt}{118}: a no-op unless the option is on, it adds the compile flags
(\code{-O0 -g -fprofile-instr-generate -fcoverage-mapping} \src{CMakeLists.txt}{108})
and --- only when the target is not a static library --- the matching link flags
\src{CMakeLists.txt}{121}. The distinction is load-bearing: static archives never see a
link line, so the profiling runtime is linked once into each final test executable
rather than uselessly requested on every archive.

The \code{coverage} target itself \src{CMakeLists.txt}{256} runs the whole pipeline
through a script-mode CMake driver, \srcf{cmake/RunLlvmCoverage.cmake}. In order: ctest
runs the aggregate suites under
\cppinline|LLVM_PROFILE_FILE=<raw>/%p-%m.profraw|
\src{cmake/RunLlvmCoverage.cmake}{33} --- \code{\%p} (process id) and \code{\%m} (binary
signature) give each test process its own profile file instead of clobbering one ---
and a nonzero ctest result is fatal \src{cmake/RunLlvmCoverage.cmake}{39}: the pipeline
is fail-closed, refusing to report coverage of a failing tree. The raw profiles are
merged with \cppinline|llvm-profdata merge -sparse|
\src{cmake/RunLlvmCoverage.cmake}{50}; \code{llvm-cov report} writes the text summary to
\code{build-coverage/coverage/summary.txt} \src{cmake/RunLlvmCoverage.cmake}{80}; and
\code{llvm-cov show} emits the HTML report with per-region line counts and branch counts
to \code{build-coverage/coverage/html/} \src{cmake/RunLlvmCoverage.cmake}{98}. The
ignore regex excludes \code{\_deps}, the build trees, \srcf{tests/} itself, and the
vendored \srcf{gui/qcustomplot.cpp} \src{CMakeLists.txt}{262}, so the percentages
measure project code exercised by the suites --- not dependency code and not the tests'
own bodies.

\begin{listing}[htbp]
\begin{minted}{cmake}
set(QTROCKET_COVERAGE_BINARIES
    "$<TARGET_FILE:sim_tests>"
    "$<TARGET_FILE:model_tests>"
    "$<TARGET_FILE:integration_tests>"
    "$<TARGET_FILE:propagator_tests>"
    "$<TARGET_FILE:cli_tests>")
\end{minted}
\caption{The coverage target's binary list \srccap{CMakeLists.txt}{248} --- five
binaries, not six. The same five names appear in the target's \code{DEPENDS}
\srccap{CMakeLists.txt}{267}, while the test-selection regex matches all seven
aggregates \srccap{CMakeLists.txt}{261}.}\label{lst:testing-covbins}
\end{listing}

\begin{keyinsight}[The coverage target predates the design-matrix suite]
A verified latent defect: \code{design\_matrix\_tests} is missing from both the
\code{-object} list and the \code{DEPENDS} of the \code{coverage} target
(\cref{lst:testing-covbins}), yet the selection regex
\cppinline|^qtrocket_.*tests$| \src{CMakeLists.txt}{261} matches both of that binary's
aggregate entries --- the coverage stack landed in one commit (\code{2c9404b}) before the
design-matrix suite existed, and the lists were never updated. Three consequences
follow. In a fresh coverage tree, the target does not build the binary, the inner ctest
run hits an unbuilt executable, and the script aborts at its fail-closed gate
\src{cmake/RunLlvmCoverage.cmake}{39} --- unless the user has already built the whole
tree. When the binary is built, every coverage run pays for the heavy full-ladder sweep,
which is slow by design \src{tests/CMakeLists.txt}{85}. And although the binary is
instrumented \src{tests/CMakeLists.txt}{72} and its profiles are merged, it is never
passed to \code{llvm-cov} as an \code{-object}, so code exercised \emph{only} by it
would be invisible in the report --- benign today because it links the same \code{cli}
and \code{qtrocket\_core} archives whose mappings arrive through the other binaries, and
its own translation unit sits under the ignored \srcf{tests/} path.
\end{keyinsight}

\subsection{Continuous integration}\label{sec:testing-ci}

CI is one GitHub Actions workflow, triggered on pushes and pull requests against
\code{development} \src{.github/workflows/cmake-multi-platform.yml}{12}. Its design
statement leads the file: CI drives the build through the project's presets, ``so it
exercises the same Ninja-based configure/build/test path that developers use locally
instead of each platform's default generator''
\src{.github/workflows/cmake-multi-platform.yml}{3}. There is no separate CI build
system to drift out of sync --- \cref{tab:testing-ci} is simply the preset table of
\cref{sec:overview} mapped onto runners.

\begin{table}[htbp]
\centering\small
\begin{tabular}{@{}lll@{}}
\toprule
Runner & Preset & Toolchain \\
\midrule
\code{ubuntu-latest} & \code{release-gcc} & GCC \\
\code{ubuntu-latest} & \code{release-clang} & Clang \\
\code{macos-latest} & \code{release-clang} & AppleClang \\
\code{windows-latest} & \code{release-msvc} & MSVC (Ninja + \code{cl}) \\
\midrule
\code{ubuntu-latest} (QEMU VM) & \code{release-clang} & FreeBSD 15.0, Clang/libc++ \\
\bottomrule
\end{tabular}
\caption{The CI matrix \srccap{.github/workflows/cmake-multi-platform.yml}{28}: four
hosted-runner entries with \code{fail-fast} disabled --- so one platform's breakage does
not mask the others --- plus the separate FreeBSD virtual-machine job (lines 26 and 94
of the same workflow).}\label{tab:testing-ci}
\end{table}

The matrix steps are few: install CMake and Ninja (Ninja is not preinstalled on the
Linux/macOS images, and every preset requires it)
\src{.github/workflows/cmake-multi-platform.yml}{37}; install Qt 6.10.2
\src{.github/workflows/cmake-multi-platform.yml}{42}; on Windows, set up the MSVC
toolchain environment explicitly, because the Ninja generator does not run vcvars the
way the Visual Studio generator implicitly did
\src{.github/workflows/cmake-multi-platform.yml}{48}; warm a ccache keyed on
OS $\times$ preset, explicitly to cache the compiled FetchContent dependencies across
runs \src{.github/workflows/cmake-multi-platform.yml}{54}; then configure, build, and
test through the presets \src{.github/workflows/cmake-multi-platform.yml}{61}. The test
step's filter is the aggregate contract of \cref{sec:testing-binaries} --- and because
\emph{both} design-matrix entries match it, every push and pull request runs the heavy
full-ladder sweep in CI \src{.github/workflows/cmake-multi-platform.yml}{74}. (The
comment above that step still names only four suites --- ``(model, sim, integration,
propagator)'' \src{.github/workflows/cmake-multi-platform.yml}{73}; it predates the CLI
and design-matrix suites. Read the filter, not the comment.)

The fifth platform has no hosted runner, so the workflow boots a FreeBSD 15.0 virtual
machine (QEMU, via \code{vmactions/freebsd-vm}) inside an \code{ubuntu-latest} runner
\src{.github/workflows/cmake-multi-platform.yml}{76}. The job is separate from the
matrix because the toolchain setup differs: Qt comes from \code{pkg} and the base
compiler is Clang with libc++, but the build itself reuses the \code{release-clang}
preset unchanged --- \code{pkg install -y cmake ninja git pkgconf qt6-base qt6-tools}, then
the same three preset commands under \code{set -e}
\src{.github/workflows/cmake-multi-platform.yml}{101}. The hermetic curl configuration
of \cref{sec:overview} pays off exactly here, and the workflow says so: ``The hermetic
curl config means OpenSSL + zlib (both in the FreeBSD base system) are the only external
libs curl needs'' \src{.github/workflows/cmake-multi-platform.yml}{99}. A 90-minute
timeout bounds the run, since the VM has no ccache and compiles every FetchContent
dependency from source \src{.github/workflows/cmake-multi-platform.yml}{87}.

What CI does not run, verified by the single workflow file being the only one under
\srcf{.github/workflows/}: no coverage job, no sanitizer job, no debug-preset job, no
artifact upload, and no release or packaging step. \Cref{sec:incomplete} carries the
inventory.

\subsection{Warning hygiene}\label{sec:testing-warnings}

\Cref{sec:overview} stated the policy --- \code{-Wall -Wextra -Wpedantic -Werror} on GCC
and Clang \src{CMakeLists.txt}{97}, \code{/W4 /WX} on MSVC \src{CMakeLists.txt}{88} ---
and the detail that makes it livable is placement. The
\cppinline|add_compile_options| calls are directory-scoped and sit \emph{after} all five
\code{FetchContent\_MakeAvailable} blocks but \emph{before} the project's own
subdirectories and targets, so vendored dependencies compile with their own flags while
every project target must be warning-clean; the history note records that the exemption
was verified in \code{compile\_commands.json} when the flags landed \src{TODO.md}{23}.
One vendored file lives inside a project target and is exempted per source: the
third-party plotting widget \srcf{gui/qcustomplot.cpp} is compiled with \code{-w} ---
``qcustomplot throws some warnings, so silence them so we can turn on -Werror for our
code'' \src{CMakeLists.txt}{152}.

On MSVC, \code{/WX} turns the compiler's C4996 ``use the \code{\_s} variant''
deprecation nags into hard errors, so the build defines
\code{\_CRT\_SECURE\_NO\_WARNINGS} tree-wide, with the rationale recorded in place:
those deprecations are non-standard editorializing against standard C++ functions like
\code{std::getenv}, genuine \cppinline|[[deprecated]]| uses still warn, and GCC/Clang
remain the real safety net \src{CMakeLists.txt}{89}. One line of that comment block has
gone stale and should not be trusted: ``We also run ASAN and clang-tidy for additional
security'' \src{CMakeLists.txt}{93}. No \code{-fsanitize} flag, no \code{CLANG\_TIDY}
property, and no \code{.clang-tidy} configuration exists anywhere in the project's
CMake, presets, or CI --- each a failed search of the tree. The warning flags are real
and enforced; the sanitizer claim is a comment about infrastructure that does not
exist.

\subsection{What the infrastructure does not do}\label{sec:testing-absent}

Distribution is where the infrastructure stops. The project contains exactly one
\cppinline|install()| rule, for the GUI executable alone ---
\cppinline|install(TARGETS qtrocket BUNDLE DESTINATION . LIBRARY DESTINATION ...)|
\src{CMakeLists.txt}{308}. Nothing installs \code{qtrocket-cli},
\code{qtrocket-visualizer}, the motor data, or any header, and the one rule works on
Linux only through side effects: it names no \code{RUNTIME DESTINATION}, so the
executable lands in \code{bin/} by CMake's implicit default, and
\code{CMAKE\_INSTALL\_LIBDIR} is defined only because some dependency's build includes
\code{GNUInstallDirs} --- no project CMake file includes it (a failed search). There is
no CPack configuration anywhere in the project's CMake files, hence no installers or
packages of any kind.

\begin{keyinsight}[\code{cmake -{}-install} would ship the vendored dependencies]
None of the five \code{FetchContent\_Declare} blocks passes \code{EXCLUDE\_FROM\_ALL}
\src{CMakeLists.txt}{15}, so the generated root install script includes the
\code{\_deps} sub-installs for GoogleTest, jsoncpp, curl, Eigen, and Boost alongside the
project's own rule --- observable in the generated \code{build/cmake\_install.cmake}.
Running \code{cmake -{}-install build} today would install GoogleTest, curl, and Boost
headers and archives into the prefix next to \code{qtrocket}. The install rule is a
template stub, not a distribution story; treat it accordingly.
\end{keyinsight}

Internationalization is dead scaffolding from the same template: the build requires
\code{LinguistTools} \src{CMakeLists.txt}{139}, lists the single translation source
\code{qtrocket\_en\_US.ts} \src{CMakeLists.txt}{145} among the GUI sources
\src{CMakeLists.txt}{187}, and runs \cppinline|qt_create_translation(QM_FILES ...)|
\src{CMakeLists.txt}{201} --- but \code{QM\_FILES} appears nowhere else in any CMake
file, so the compiled catalogs are never attached to a target or installed; the roadmap
parks translations explicitly \src{TODO.md}{79}. The macOS bundle identifier is likewise
still the template placeholder \code{my.example.com} \src{CMakeLists.txt}{301}.

The shape of the whole is easy to state. Everything between a developer's editor and a
green five-platform CI run is built, pinned, and self-checking: six suites whose
assertions are oracles, a fixture corpus that audits its own manifest, a coverage
pipeline that refuses to report on a failing tree, and a warning wall that vendored code
cannot breach. Everything between that green run and a user's machine --- installers,
packages, translations, a runtime motor database --- is not built yet, and
\cref{sec:incomplete} inventories it with the rest of the unfinished work.
